\documentclass[12pt]{article}
\usepackage[utf8]{inputenc}
\usepackage[T1]{fontenc}
\usepackage{hyperref}
\title{Telugu Language Resources for K-12 Curriculum Development}
\author{Compiled by Claude Sonnet 5 (claude-sonnet-5) for Andr\'as Kornai}
\date{}
\begin{document}
\maketitle
\section{Introduction}
Telugu (ISO 639-3 \texttt{tel}) is a Dravidian language and a classical language of India, the most widely spoken Dravidian language, and a language of instruction in Telangana and Andhra Pradesh. This report is limited: only a few searches were run, and many hits were unreliable aggregator sites, which are omitted.
\section{Primary Educational Resources}
\begin{itemize}
\item Telangana SCERT e-Textbooks portal (textbooks for classes I--X, including Telugu medium; reported by Careers360, the portal URL itself was not shown as a search result, so verify): see \url{https://school.careers360.com/boards/bse-telangana/telangana-board-books-for-class-10-bsar}
\end{itemize}
No official Andhra Pradesh SCERT page was confirmed in search results.
\section{Linguistic and Reference Materials}
\begin{itemize}
\item Wikipedia, Telugu language: \url{https://en.wikipedia.org/wiki/Telugu_language}
\item Wikipedia, Charles Philip Brown (author of the Telugu-English dictionary, 1852/1854): \url{https://en.wikipedia.org/wiki/Charles_Philip_Brown}
\item Lexilogos Telugu dictionary links: \url{https://www.lexilogos.com/english/telugu_dictionary.htm}
\end{itemize}
\section{Religious and Community Resources}
None found in these searches.
\section{Digital and Audio Resources}
\begin{itemize}
\item Lexilogos (online dictionary portal, listed above).
\end{itemize}
\section{Curriculum Development Context}
Telugu has an extensive state-board textbook system. The next step is to inventory official SCERT textbooks and licensing for each grade, and then identify gaps (e.g.\ digital audio, assessment materials). Telugu Academy and other publishers were not confirmed in these searches.
\end{document}